\documentclass{beamer}
\usepackage{bbm}
\usepackage{booktabs}
\usepackage{array}
\usepackage{bm}
\setbeamertemplate{footline}[frame number]


%Information to be included in the title page:
\title{Grid-Scale Battery Dispatch Distortions \\ Due to Manufacturer Warranties}
\author{Julia Park}
\date{IO Workshop \\ April 20, 2026}

\begin{document}

\frame{\titlepage}

\begin{frame}{Motivation}
\begin{itemize}
    \item Grid-scale batteries have expanded rapidly in U.S. electricity markets
    \begin{itemize}
        \item Over 26 GW of installed capacity by early 2025
        \item Texas (ERCOT) alone hosts approximately 3.2 GW (Reynolds, 2025)
    \end{itemize}
    \pause
    \item A primary economic rationale for this investment is \textbf{flexibility}
    \begin{itemize}
        \item Unlike intermittent renewables or fossil-fuel generators with startup costs, batteries can respond within seconds
        \item Uniquely valuable for smoothing unexpected spikes in demand
    \end{itemize}
    \pause
    \item But in practice, battery operation is not unconstrained
    \begin{itemize}
        \item Batteries face contractual operational restrictions from offtake agreements / PPAs, warranties, etc.
    \end{itemize}
\end{itemize}

\textbf{Research question}: \textit{How do contractual operational restrictions for grid-scale battery distort dispatch away from profit-maximizing or socially optimal behavior?}
\end{frame}


\begin{frame}{Roadmap}
    \begin{enumerate}
        \item Institutional background on battery operational restrictions
        \item Literature review
        \item Data sources
        \item Proposed bidding model simulation
        \item Future steps
    \end{enumerate}
\end{frame}

\section{Institutional Background}

\begin{frame}{Context: Where do operational restrictions come from?}
\begin{itemize}
    \item \textbf{Performance Warranties}: guarantee from battery manufacturer that the storage system will maintain specific performance standards over 15-20 years, if battery is operated within certain constraints
    \item \textbf{Power Purchase Agreements (PPAs) / Energy Storage Agreement (ESAs)}: financial agreement to purchase power between operator and offtaker. 
    \begin{itemize}
        \item These days, the operational restrictions from these are usually harmonized (although interviews have told me that in the past that wasn't always the case...)
    \end{itemize}
\end{itemize}
\end{frame}



\begin{frame}{Context: Common types of operating restrictions}
Restrictions imposed on battery operation to prevent degradation:
\begin{itemize}
    \item \textbf{Cycle frequency limits}: cap on full charge-discharge cycles per day or per year (typically clustered around 1--1.5 cycles/day)
    \item \textbf{State-of-charge requirements}: constrain annual average SOC, commonly below or above 50\%
\end{itemize}
\end{frame}

\begin{frame}{Context: Trends in battery warranties}
The industry has been interested in how to design these operational restrictions such that they are less \textbf{restrictive}  (e.g. limits on number of cycles) and more \textbf{flexible} warranties (e.g. rolling averages, charge/discharge rates over the day).

\vspace{0.3cm}
Reasons for change:
\begin{itemize}
    \item Batteries earn much of their money arbitraging on a few high-demand days, rather than just once every day, so restrictions can prevent optimal dispatch
    \begin{itemize}
        \item \textit{``Running a BESS with an inflexible guarantee is like paying rent all year for an umbrella shop, but being forced to keep the same daily stock; so when the rainy week comes, you sell out in the first hour and miss the payoff"} \textcolor{gray}{-Gregory Green, BlackVolt}
    \end{itemize}
\end{itemize}
\end{frame}

\section{Literature Review and Positioning}

\begin{frame}{Literature: Electricity Storage Arbitrage}
How storage facilities optimize dispatch, how they affect market prices, and frictions preventing storage from realizing theoretical value.

\vspace{0.3cm}
\begin{itemize}
    \item \textbf{Butters, Dorsey, and Gowrisankaran (2025)}: dynamic market equilibrium model of battery dispatch in CAISO; first 5,000 MWh of storage reduces wholesale prices by 5.6\%
    \begin{itemize}
        \item Their framework (dynamic programming with physical degradation costs) provides a starting point for my bidding simulation
        \item But their degradation costs are physical---warranty constraints may be substantially more restrictive
    \end{itemize}
    \item \textbf{Lamp and Samano (2022)}: document CAISO battery dispatch 2018--2019, find behavior only partially consistent with optimal arbitrage
    \begin{itemize}
        \item Attribute deviations to forecasting and strategic behavior
        \item But they don't observe warranty terms---my project asks whether warranty terms explain part of the residual
    \end{itemize}
\end{itemize}
\end{frame}

\begin{frame}{Literature: Electricity Storage Arbitrage, cont'd}
\begin{itemize}
    \item \textbf{Reynolds (2025)}: documents that 66\% of ERCOT storage is classified for arbitrage and 44\% for ancillary services (with overlap)
    \begin{itemize}
        \item Context for welfare that batteries provide
    \end{itemize}
    \item \textbf{Denholm and Margolis (2018)}: estimates storage required to meet California peaking demand under increasing solar
    \begin{itemize}
        \item Useful for situating the social stakes of dispatch inefficiency
    \end{itemize}
\end{itemize}
\end{frame}

\begin{frame}{Literature: Contract Theory and Warranties}
How warranty terms emerge from bilateral contracting, how they allocate risk, govern moral hazard.

\vspace{0.2cm}
Prior focus: consumer durable goods---cars, appliances, electronics---where buyer usage affects product longevity.

\vspace{0.2cm}
The battery setting differs in two ways:
\begin{enumerate}
    \item The buyer's dispatch decisions affect \textit{wholesale electricity market outcomes}, so warranty design has social welfare consequences beyond the bilateral relationship
    \item Unlike most consumer goods (degrading along 1--2 dimensions), batteries degrade via \textit{many} behaviors (DoD vs. cycles vs. SOC vs. C-rate, etc.), and warranty terms vary in restrictiveness across these behaviors---often depending on monitorability
\end{enumerate}
\end{frame}

\begin{frame}{Literature: Contract Theory, cont'd}
\begin{itemize}
    \item \textbf{Emons (1988)}: under asymmetric information, warranties serve as both insurance and quality signals; moral hazard induces excessive product use
    \begin{itemize}
        \item Warranty cycling restrictions are one instrument for managing this
    \end{itemize}
    \item \textbf{Emons (1989)}: longer warranties require more intensive monitoring to remain incentive-compatible
    \begin{itemize}
        \item Directly relevant to 15--20 year battery warranties, where cumulative cycle counts are both technically feasible and contractually specified
    \end{itemize}
    \item \textbf{Liu and Wang (2025)}: industry review of potentially optimal battery warranty designs from the manufacturer's perspective
    \begin{itemize}
        \item This project complements theirs by examining the \textit{market-level} implications
    \end{itemize}
\end{itemize}
\end{frame}

\begin{frame}{Literature: Battery Degradation}
Engineering foundation for the economic analysis.

\vspace{0.2cm}
\begin{itemize}
    \item \textbf{Xu et al. (2018)}: most widely used semi-empirical degradation model
    \begin{itemize}
        \item Separates aging into calendar aging (temperature, SOC level) and cycle aging (DoD, C-rate, cumulative throughput)
    \end{itemize}
    \item \textbf{Attia et al. (2022)}: document non-linear aging in lithium-ion batteries, including ``capacity knees''---abrupt inflection points where degradation accelerates
    \begin{itemize}
        \item Simple linear throughput approximations may misstate marginal degradation costs
        \item Supports conservative warranty design near the knee threshold
        \item Implies batteries operating well below the knee face lower marginal degradation costs than linear models suggest
    \end{itemize}
\end{itemize}
\end{frame}

\begin{frame}{My Contribution}
    \begin{itemize}
        \item Unlike consumer products studied in the warranty literature, grid-scale batteries have \textbf{multiple contractable dimensions of operational behavior}---cycle frequency, SOC
        \item These choices have potential consequences for \textit{electricity market outcomes} and welfare
    \end{itemize}
\begin{itemize}
        \item \textbf{Document and categorize} operational restrictions embedded in performance warranties using PPA-disclosed warranty terms
    \item \textbf{Simulate dispatch under different regimes}, find costs (if any) of different contract designs
\end{itemize}

\end{frame}



\section{Part 1: Documenting Battery Operational Restrictions}


\begin{frame}{Part 1: Documenting Battery Operational Restrictions}
\begin{itemize}
    \item Utilities in Texas (PUCT) and New Mexico (NMPRC) must file PPAs and Energy Storage Agreements (ESAs) for regulatory approval, which sometimes specify exact cycle limits, SOC windows, and enforcement rules
    \item (In the process of scraping) operational restrictions from publicly released ESAs for battery projects in New Mexico and Texas
    \begin{itemize}
        \item \textit{Why New Mexico?} tends to most often have unredacted ESA contract language
        \item \textit{Why Texas?} Any ESA can be linked to granular public data on the individual battery's individual bids and dispatch behavior -- can be useful later on...
    \end{itemize}
    \item So far managed to scrape restrictions from \textbf{11 dockets}, covering \textbf{16 battery projects} filed between 2021 and 2025
\end{itemize}
\end{frame}

\begin{frame}{Texas: Operational Restrictions Across Dockets}
\footnotesize
\begin{tabular}{p{3.5cm} p{0.9cm} p{2.0cm} p{2.5cm}}
\toprule
\textbf{Project} & \textbf{MW} & \textbf{Annual cycles} & \textbf{SOC constraint} \\
\midrule
Wildcat Ranch BESS (TX-57244) & 48 & 250 cap; excess to ${\sim}$365 at add'l charge & $<$50\% avg ceiling \\
\addlinespace
Carne ESA (TX-57760) & 65 & 340 cap; 2/day & $<$60\% avg ceiling; $<$20\% at rest \\
\addlinespace
Lorenzo/Chaves/ Roswell BESSes (TX-59034) & 50 ea. & 365; 5\%/day & $>$51\% avg \textit{floor}\\
\addlinespace
Wildcat Ranch II / Palo Duro / Mammoth Plains (TX-59034) & 120--150 & 365; 5\%/day & $>$51\% avg \textit{floor} \\
\bottomrule
\end{tabular}


\end{frame}

\begin{frame}{New Mexico: Operational Restrictions Across Dockets}
\footnotesize
\begin{tabular}{p{3.5cm} p{0.8cm} p{1.9cm} p{3cm}}
\toprule
\textbf{Project} & \textbf{MW} & \textbf{Annual cycles} & \textbf{SOC constraint} \\
\midrule
Sky Ranch (NM-21-00031) & 100 & 365; optional to 535 & $<$40\% BAASOC ceiling; penalty: $-$5 cycles  \\
\addlinespace
TAG Energy (NM-23-00251) & 50 & 365 hard cap & Manufacturer bounds only; no avg ceiling  \\
\addlinespace
Quail Ranch (NM-23-00353) & 100 & 365; 2/day & $<$60\% avg ceiling; $>$30\% at rest ($>$3 mo.)  \\
\addlinespace
Route 66 / Sky Ranch II (NM-23-00353) & 50/100 & 365; daily = mfr warranty only & $<$40\% BAASOC ceiling  \\
\addlinespace
Four Mile Mesa / Star Light / Windy Lane (NM-25-00048) & 68--100 & 365; 2/day & $<$50\% avg ceiling; min 3\% for $\leq$1 hr \\
\bottomrule
\end{tabular}
\end{frame}

\begin{frame}{Part 2: Simulating optimal dispatch}
Simulate optimal dispatch under five regimes, given the \textit{same} expectations around future prices and physical degradation costs:
\vspace{0.2cm}
\begin{itemize}
    \item \textbf{Unconstrained}: physical degradation costs only (calibrated from Xu et al., 2018)
    \item \textbf{2-cycle-per-day limit}
    \item \textbf{Annual throughput cap} equivalent to 365 cycles
    \item \textbf{SOC ceiling ($<50\%$)}
    \item \textbf{SOC floor ($>51\%$)}
\end{itemize}

\vspace{0.3cm}
Produces \textit{upper-bound} estimate of profit loss attributable to each type of restriction

\begin{itemize}
    \item In the future: compare to socially-optimal dispatch? (peak load reduction, fossil emissions)
\end{itemize}

\end{frame}




\begin{frame}{A very simple dispatch model...}

\begin{itemize}
    \item At the start of hour $h$, the battery commits a dispatch quantity (MW) based on their expectation of future price $p_h$
    \item The committed quantity clears at the realized $p_h$
    \item These batteries are basically self-scheduled batteries in ERCOT
\end{itemize}
\end{frame}

\begin{frame}{Model setup}
\textbf{Choice and state variables:}
\begin{itemize}
    \item $g_h \geq 0$: discharge quantity; $\;l_h \geq 0$: charge quantity; $\;g_h \cdot l_h = 0$
    \item $s_h \in [0, k]$: state of charge 
    \item $p_h$: settlement-point price
\end{itemize}

\vspace{0.2cm}
\textbf{SoC dynamics} with round-trip efficiency $\delta = 0.85$ (guaranteed by warranties):
\[
s_{h+1} = s_h - g_h + \delta\, l_h
\]

\textbf{Per-period profit} with per-MWh cycling cost $c/k$:
\[
\pi_h = (g_h - l_h)\,p_h \;-\; \frac{c\,g_h}{k}
\]
\begin{itemize}
    \item Cycling cost: calibrate to physical degradation experiments
\end{itemize}

\end{frame}

\begin{frame}{Model setup}


\textbf{AR(1) price beliefs:}
\[
\bar p_h \;\equiv\; \mathbb{E}[p_h\mid p_{h-1}] \;=\; \theta_h\, (p_{h-1} - p_{\text{DAM}, h-1}) + \phi_h p_{\text{DAM}} + \mu_h.
\], iid normal errors.

\vspace{0.2cm}
\textbf{Bellman equation}: 
\begin{align*}
    V_h(s_h, p_{h-1}) =  \max_{g_h,\,l_h}\Big\{(&g_h - l_h)\,\bar p_h - \tfrac{c\,g_h}{k} + \\ &\mathbb{E}[ V_{h+1}(s_h - g_h + \delta l_h,\, p_{h-1}) | p_{h-1}]\Big\}
\end{align*}

\end{frame}

\begin{frame}{Future steps / requested feedback}
    \begin{itemize}
        \item Continuing to scrape more PPAs/ESAs for operational restrictions
        \begin{itemize}
            \item Acquiring more flexible/complicated operational terms from energy solution companies
        \end{itemize}
        \item Adding various different constraints/regimes to dispatch simulation model
        \item How to pitch the economics of this research question?
    \end{itemize}
\end{frame}


\end{document}
